\section{Related work}
\label{sec:related}

\paragraph{Fragility of multi-task evaluation.}
Benchmark outcomes depend on which tasks are included and how they are aggregated: \citet{dehghani2021benchmark} show that the relative standing of methods can change when the benchmark tasks are changed, and \citet{lanctot2023evaluating} view each task as a voter and an aggregate as a social welfare function, comparing rules through their axiomatic properties and empirical robustness. Large-model leaderboards use plain averages: HELM \citep{liang2022holistic} reports, among others, accuracy per scenario, and the Open LLM Leaderboard reports the mean of per-task normalised scores, computed with the EleutherAI evaluation harness \citep{gao2023harness}. These works vary the \emph{benchmark} (tasks, rule) for fixed models. We instead fix the benchmark and ask what a single \emph{model's developer} must change to win it, and how that cost depends on the rule.

\paragraph{Gaming of leaderboards and contamination.}
Training on test-related data is a documented route to inflated scores \citep{zhou2023cheater}; rephrased test samples evade common decontamination filters and let a small model match much larger ones on MMLU and other benchmarks \citep{yang2023rephrased}. \citet{singh2025illusion} document that undisclosed private testing and unequal sampling distort the Chatbot Arena ranking. These studies show that targeted training on a benchmark's \emph{data} moves scores. Our question is the next one: given that scores on selected sub-datasets can be raised, \emph{how much raising is needed to reach first place}, and under which aggregation rule is that cheapest or impossible.

\paragraph{Strategic post-training.}
\citet{chen2026incentives} model a benchmark designer and several developers as a Stackelberg game in which developers allocate resources to benchmark-specific improvements of a latent-quality model, and prove results on equilibria of the induced games. Their object is the equilibrium of strategic developers; ours is the cost of one developer's best response against a static field, measured exactly on observed score matrices. The two are complementary: we do not model rivals' reactions (Section~\ref{sec:limits}).

\paragraph{Social choice.}
Viewing models as candidates and sub-datasets as voters places the question in computational social choice \citep{brandt2016handbook}. In that literature, strategic action is usually on the voter side: manipulation of the ballots \citep{bartholdi1989computational} and bribery, i.e.\ paying voters to change their votes so that a chosen candidate wins \citep{faliszewski2009bribery}; the cost is the number of voters touched. Here the action is on the candidate side: the target's own task scores change, voters' preferences move as a consequence, and the cost is measured in score points and in tasks touched. Cardinal rules (means, median, minimum) and ordinal rules (Borda, plurality, Condorcet-type) are thus both natural to compare.
